\documentclass[10pt,a4paper]{amsart}
\usepackage[margin=19mm]{geometry}
\usepackage{amsmath,amssymb,mathtools}
\usepackage[scr=rsfs,cal=euler]{mathalfa}
\usepackage{xfrac}
\usepackage[unicode,hidelinks,bookmarks=false]{hyperref}
\newcommand{\ie}{i.e.\ }
\newcommand{\cf}{cf.\ }
\newcommand{\resp}{respectively\ }
\newcommand{\Subsection}{\subsection}
\providecommand{\nobreakdash}{\nobreak\hbox{-}}
\setlength{\emergencystretch}{2em}
\multlinegap0pt
\usepackage{amssymb}
\usepackage{mathtools}
\usepackage{bm}
\usepackage[shortlabels]{enumitem}
\usepackage{xcolor}
\newtheorem{lemma}{Lemma}
\newcommand{\RR}{\mathbb{R}}
\newcommand{\calM}{\mathcal{M}}
\newcommand{\calS}{\mathcal{S}}
\newcommand{\calT}{\mathcal{T}}
\newcommand{\norm}[1]{\left\lVert #1\right\rVert}
\title{Analytic inverse problems with\\ finitely many random measurements}
\author{Giovanni S. Alberti, Damiano Poletti, Simone Sanna, Matteo Santacesaria}
\date{Source: arXiv:2608.14324v1 (CC BY 4.0)}
\begin{document}
\maketitle

\noindent\emph{Source and licence.} Analytic inverse problems with\\ finitely many random measurements. Version arXiv:2608.14324v1 (CC BY 4.0), distributed by arXiv under the Creative Commons Attribution 4.0 International licence, http://creativecommons.org/licenses/by/4.0/, as recorded in the arXiv licence metadata for this exact version and shown on the abstract page https://arxiv.org/abs/2608.14324v1. Full licence notice: \texttt{licenses/arxiv-2608.14324-CC-BY-4.0.txt}.\par\smallskip

\noindent\textit{Editorial scope.} Counted unit: the article's lemma eq:secant-lower-bound (source0.tex lines 865--882). The article's complete original proof of it is delivered with it (source0.tex lines 884--909, one \texttt{proof} environment of 26 lines). No mathematics is written, repaired or re-derived: the statement and the proof are the article's own lines, in the article's own notation, and no retained line is substituted or re-flowed.  No further source-local context is needed: every cross-reference of every retained line resolves inside the two delivered spans.  Nothing was omitted from any retained span, so the package carries no omission record.  No retained line contains a citation, so the wrapper carries no bibliography and no bibliography entry.  No retained line contains a figure environment, an image-inclusion command or drawing code, so no image is delivered and the package declares no asset dependency.  Editorial formatting only, in addition: an AMS \texttt{amsart} preamble that restates the article's own declarations and macros used by the retained lines, namely the environments \texttt{lemma} on separate counters, together with the standard packages those lines need.  The retained environments print in this document as Lemma 1 in order of appearance, whereas the article prints the same items with the numbering of its own full document.  The complete native source of the article as distributed in the arXiv e-print for this exact version is bundled unchanged as \texttt{sources/207657/source0.tex}.
\begin{lemma}\label[lemma]{lem:compact-secants}
Let $\calM\subset X$ be a finite-dimensional $C^1$ embedded manifold, let $U\subset\calM$ be open, let $K\subset U$ be compact, and let $G\in C^1(U;Y)$, where $Y$ is a Banach space. Then the normalized secant set
\[
    \calS_K(G)
    =\left\{\frac{G(x)-G(y)}{\norm{x-y}_X}:x,y\in K,\ x\neq y\right\}
\]
is relatively compact in $Y$. Moreover,
\[
    \calT_K(G)
    =\left\{dG_x(v):x\in K,\ v\in T_x\calM,\ \norm{v}_X=1\right\}
\]
is compact.

If $G$ is injective on $K$ and $dG_x$ is injective for every $x\in K$, then
\begin{equation}\label{eq:secant-lower-bound}
    0\notin\overline{\calS_K(G)}\cup\calT_K(G).
\end{equation}
\end{lemma}

\begin{proof}
Take a sequence of normalized secants $\frac{G(x_j)-G(y_j)}{\norm{x_j-y_j}_X}$ for $x_j,y_j\in K$. After passing to a subsequence, the $x_j,y_j$ converge to $x,y\in K$ respectively. If $x\neq y$, convergence is immediate. If $x=y$, choose a $C^1$ chart $\varphi:V\subset\RR^d\to\calM$ around $x$, with $V$ convex after shrinking. Write the endpoints as $x_j=\varphi(p_j)$ and $y_j=\varphi(q_j)$, where $p_j,q_j\to q=\varphi^{-1}(x)$, and pass to a subsequence such that
\[
    \frac{p_j-q_j}{|p_j-q_j|}\longrightarrow v,\qquad |v|=1.
\]
The fundamental theorem of calculus gives
\[
    \frac{G(\varphi(p_j))-G(\varphi(q_j))}{|p_j-q_j|}
    \longrightarrow d(G\circ\varphi)_{q}(v),
\]
and
\[
    \frac{\norm{\varphi(p_j)-\varphi(q_j)}_X}{|p_j-q_j|}
    \longrightarrow \norm{d\varphi_q(v)}_X>0.
\]
Consequently,
\[
    \frac{G(\varphi(p_j))-G(\varphi(q_j))}
         {\norm{\varphi(p_j)-\varphi(q_j)}_X}
    \longrightarrow
    \frac{d(G\circ\varphi)_q(v)}{\norm{d\varphi_q(v)}_X},
\]
so the normalized secants have a convergent subsequence. Hence, $\mathcal S_K(G)$ is relatively compact in $Y$. The unit tangent bundle over $K$, with the norm induced by $X$, is compact by local triviality and a finite chart cover of $K$; continuity of $dG$ therefore gives compactness of $\calT_K(G)$.

If a limit of secants with distinct limiting endpoints were zero, injectivity of $G$ on $K$ would fail. If the endpoints coalesce, the preceding calculation shows that the limit is a nonzero normalized derivative whenever $dG_x$ is injective. The same injectivity excludes zero from $\calT_K(G)$.
\end{proof}

\end{document}
